% Single-column supplementary material: application-case experiments (after bibliography).
\clearpage
\onecolumn
\FloatBarrier
\begin{center}
{\Large\bfseries Supplementary Material}\\[0.4em]
{\normalsize Application-Case Experiments (MACPO Supplementary Parts V--VI and IEEE Pilots)}
\end{center}
\vspace{0.6em}
\FloatBarrier
\setcounter{section}{4}
\renewcommand{\thesection}{\Roman{section}}

\section{Experiment on the Resource Allocation Problem}\label{app:resource}
Following MACPO Supplementary Part~V~\cite{ref_macpo}, each node~$i$ produces resource $g_i$ and exchanges signed transfers $c_{i,j}$ with neighbors $j\in\mathcal{N}_i$:
\begin{equation}
\min F(\mathbf{X})=\sum_{i=1}^{n}\Bigl(U_i-\bigl(g_i+\textstyle\sum_{j\in\mathcal{N}_i}c_{i,j}\bigr)\Bigr)^2,
\qquad
g_i=x_{I_i},\quad
c_{i,j}=\begin{cases}
x_{I_{i,j}}, & i<j,\\
-x_{I_{i,j}}, & i>j,
\end{cases}
\label{eq:app_resource}
\end{equation}
where $U_i$ is the expected resource level at node~$i$ and $x_{I_{i,j}}$ is the shared variable on edge $\{i,j\}$. The objective penalizes the gap between expected and realized supply after production and transfers.

\paragraph{Experimental setup.}
The resource-allocation pilot uses $n{=}10$ agents on a chain graph ($D{=}19$). Demands are $U_i\in\{9.0,\ldots,12.5\}$; decisions lie in $[-5,25]$. Each agent evaluates Eq.~\eqref{eq:app_resource} locally. The evaluation cap is $6.0\times10^4$ function evaluations; each agent runs LLSO with swarm size 100 and penalty weight $w=|f|/512$.

\section{Experiments on an Electric Dispatching Problem}\label{app:ev}
Following Supplementary Part~VI~\cite{ref_macpo}, each power station produces electricity, trades along edges with loss, and sells surplus to electric vehicles (EVs):
\begin{equation}
\min F(\mathbf{X})=\sum_{i=1}^{n}\Bigl[f_{\mathrm{product}}(g_i)+\textstyle\sum_{j\in\mathcal{N}_i}f_{\mathrm{deal}}(c_{i,j})-f_{\mathrm{gain}}\!\Bigl(g_i+\textstyle\sum_{j\in\mathcal{N}_i}c_{i,j}\Bigr)\Bigr],
\label{eq:app_ev_global}
\end{equation}
with $g_i=x_{I_i}$ and the same signed $c_{i,j}$ convention as Eq.~\eqref{eq:app_resource}, and
\begin{equation}
f_{\mathrm{product}}(g_i)=a_i g_i,\qquad
f_{\mathrm{deal}}(c_{i,j})=p_1(1+\alpha L_{i,j})\,c_{i,j}.
\label{eq:app_ev_costs}
\end{equation}
EV revenue $f_{\mathrm{gain}}$ is computed by evolutionary game (Supplementary Algorithm~S1). Station and vehicle actions are
\begin{equation}
p_{\mathrm{sell}}^{t}\!=\!\begin{cases}
p_{\mathrm{sell}}^{t-1}+\epsilon, & \sum_k Q_k > x_{\mathrm{total}},\\
p_{\mathrm{sell}}^{t-1}-\epsilon, & \sum_k Q_k < x_{\mathrm{total}},\\
p_{\mathrm{sell}}^{t-1}, & \text{otherwise},
\end{cases}
\qquad
Q_k=\arg\min_{Q_k}\bigl[p_2 Q_k+R(Q_k+Q_k^0)-R(Q_k^0)\bigr],
\label{eq:app_ev_game}
\end{equation}
with $R(Q)=Q^2-(m+1)Q+m$, $m>1$, and $Q_k^0\in[0,1]$ the initial SOC fraction of vehicle~$k$.

\begin{algorithm}[H]
\caption{Gain calculation at station $i$ (Supplementary Algorithm~S1~\cite{ref_macpo})}
\label{alg:ev_gain}
\begin{algorithmic}[1]
\Require production plan $X_i$, vehicle count $N_v$, initial SOC $Q_k^0$ ($k=1,\ldots,N_v$)
\State $x_{\mathrm{total}}\gets g_i+\sum_{j\in\mathcal{N}_i}c_{i,j}$;\quad $p_{\mathrm{sell}}\gets a_i$;\quad initialize $Q_k$
\While{sell price not stable}
  \State update $p_{\mathrm{sell}}$ by Eq.~\eqref{eq:app_ev_game} (left branch)
  \For{$k=1$ to $N_v$}
    \State $Q_k\gets\arg\min_{Q_k}\bigl[p_2 Q_k+R(Q_k+Q_k^0)-R(Q_k^0)\bigr]$
  \EndFor
\EndWhile
\State $f_{\mathrm{gain}}\gets p_{\mathrm{sell}}\sum_{k=1}^{N_v} Q_k$
\end{algorithmic}
\end{algorithm}

\paragraph{Experimental setup.}
The EV-dispatch pilot uses $n{=}10$ agents, $N_v{=}5$ vehicles per station, $a_i\in[0.08,0.12]$, $p_1{=}0.05$, $\alpha{=}0.02$, $L_{i,j}{=}1$, $m{=}2.5$, $\epsilon{=}0.02$, and at most 80 price iterations per fitness evaluation. The evaluation cap and topology match the resource-allocation pilot ($6.0\times10^4$ evaluations, 10 agents, chain graph).

\section{Multi-Area Economic Dispatch (MAED-13)}\label{app:maed}
MAED-13 is a 13-unit valve-point economic dispatch (VPE-ED) benchmark at 1800\,MW total load:
\begin{equation}
C_i(P_i)=a_i P_i^2+b_i P_i+c_i+\bigl|e_i\sin\bigl(f_i(P_i^{\min}-P_i)\bigr)\bigr|.
\label{eq:app_maed}
\end{equation}
This pilot uses a single agent ($D{=}13$; no shared variables). Power balance is enforced by slack-unit repair plus a quadratic penalty; box limits use clamping and soft violations. The evaluation cap is $3.25\times10^5$ function evaluations.

\section{IEEE Regional Transmission Dispatch}\label{app:ieee}\label{app:ieee_partition}
These pilots extend the NDO template to MATPOWER IEEE~30/57/118 test cases (continuous economic dispatch, not unit commitment). Each MPI process represents a control region; private variables are generator outputs $P_g$; shared variables are tie-line powers $P_{\mathrm{tie}}$. Region~$r$ evaluates
\begin{equation}
f_r(\mathbf{x})=
\sum_{g\in\mathrm{region}\,r} C_g(P_g)
+ w_{\mathrm{bal}}\Bigl(\textstyle\sum_{g\in\mathrm{region}\,r} P_g - D_r - P_{\mathrm{export},r}\Bigr)^2
+ \frac{1}{n_{\mathrm{reg}}}\,\Pi_{\mathrm{tie}}(\mathbf{x}),
\label{eq:app_ieee_local}
\end{equation}
with $C_g(P)=a_g P^2+b_g P+c_g$, fixed regional demand $D_r$ (MW), $w_{\mathrm{bal}}{=}1000$, tie capacity penalty $\Pi_{\mathrm{tie}}$ ($w_{\mathrm{tie,cap}}{=}80$). The reported global objective is
\begin{equation}
F=\sum_{r=1}^{n_{\mathrm{reg}}} f_r(\mathbf{x})
\label{eq:app_ieee_global}
\end{equation}
at the best assembled global dispatch (rank-0 summation; does not consume the evaluation budget).

\paragraph{Variable encoding.}
Decision vector $\mathbf{x}\in\mathbb{R}^{D}$, $D=G+T$:
\begin{equation}
x_d \equiv P_{G,g}\ (d=0,\ldots,G-1),\qquad
x_{G+t}\equiv P_{\mathrm{tie},t}\ (t=0,\ldots,T-1).
\label{eq:app_ieee_encoding}
\end{equation}
Regions are obtained by multi-source breadth-first search from generator buses with load-balancing swaps. Initial tie flows are set to zero with a feasible power-balance repair.

\begin{figure}[H]
\centering
\begin{tikzpicture}[
  region/.style={draw, rounded corners=2pt, minimum width=1.8cm, minimum height=0.9cm, align=center},
  tie/.style={font=\small, fill=white, inner sep=1pt}
]
\node[region] (R0) at (0,1.2) {R0\\62.5 MW};
\node[region] (R1) at (3.0,1.2) {R1\\56.5 MW};
\node[region] (R2) at (0,-0.6) {R2\\23.7 MW};
\node[region] (R3) at (3.0,-0.6) {R3\\46.5 MW};
\draw[flow] (R0) -- node[tie, above=2pt] {$P_{\mathrm{tie},6}$} (R1);
\draw[flow] (R0) -- node[tie, left=2pt] {$P_{\mathrm{tie},7}$} (R2);
\draw[flow] (R0) to[bend left=8] node[tie, above=1pt] {$P_{\mathrm{tie},8}$} (R3);
\draw[flow] (R1) -- node[tie, right=2pt] {$P_{\mathrm{tie},9}$} (R2);
\draw[flow] (R1) -- node[tie, below=2pt] {$P_{\mathrm{tie},10}$} (R3);
\end{tikzpicture}
\caption{Four-region partition for IEEE~30-bus ED (capacities 96/97/64/64/16\,MW on ties $d=6,\ldots,10$).}
\label{fig:ieee30_partition}
\end{figure}

\begin{table}[H]
\centering
\caption{IEEE~30-bus regional partition ($G{=}6$, $T{=}5$, $D{=}11$; total demand 189.2\,MW).}
\label{tab:ieee30_partition}
\begin{tabular}{@{}clccp{0.42\textwidth}@{}}
\toprule
\textbf{Rank} & \textbf{Buses} & \textbf{Gens} & \textbf{$D_r$} & \textbf{Local dims (private + shared)} \\
\midrule
R0 & 1--7, 9--11, 20 & 1, 2 & 62.5 & $P_{G,1},P_{G,2},P_{\mathrm{tie},6},P_{\mathrm{tie},7},P_{\mathrm{tie},8}$ \\
R1 & 14--15, 18--19, 21--24 & 22, 23 & 56.5 & $P_{G,22},P_{G,23},P_{\mathrm{tie},6},P_{\mathrm{tie},9},P_{\mathrm{tie},10}$ \\
R2 & 12--13, 16--17 & 13 & 23.7 & $P_{G,13},P_{\mathrm{tie},7},P_{\mathrm{tie},9}$ \\
R3 & 8, 25--30 & 27 & 46.5 & $P_{G,27},P_{\mathrm{tie},8},P_{\mathrm{tie},10}$ \\
\bottomrule
\end{tabular}
\end{table}

\begin{table}[H]
\centering
\caption{IEEE~57-bus regional partition ($G{=}7$, $T{=}4$, $D{=}11$; total demand 1250.8\,MW).}
\label{tab:ieee57_partition}
\begin{tabular}{@{}clccp{0.38\textwidth}@{}}
\toprule
\textbf{Rank} & \textbf{Buses (subset)} & \textbf{Gens} & \textbf{$D_r$} & \textbf{Shared ties} \\
\midrule
R0 & 1--2, 14--17, 44--47 & 1, 2 & 217.2 & $P_{\mathrm{tie},7}$ (R0--R2), $P_{\mathrm{tie},8}$ (R0--R3) \\
R1 & 4--8, 18--21, 24--31, 52 & 6, 8 & 322.3 & $P_{\mathrm{tie},9},P_{\mathrm{tie},10}$ \\
R2 & 3 & 3 & 41.0 & $P_{\mathrm{tie},7},P_{\mathrm{tie},9}$ \\
R3 & 9--13, 22--23, 32--57 & 9, 12 & 670.3 & $P_{\mathrm{tie},8},P_{\mathrm{tie},10}$ \\
\bottomrule
\end{tabular}
\end{table}

\begin{table}[H]
\centering
\small
\caption{IEEE~118-bus eight-region partition ($G{=}54$, $T{=}11$, $D{=}65$; total demand 4242\,MW).}
\label{tab:ieee118_partition}
\begin{tabular}{@{}clccp{0.40\textwidth}@{}}
\toprule
\textbf{Rank} & \textbf{Buses (representative)} & \textbf{\#G} & \textbf{$D_r$} & \textbf{Shared tie dims} \\
\midrule
R0 & 1--7, 11--14, 16, 117 & 4 & 430 & 54, 55, 56 \\
R1 & 82--93, 96--97, 102 & 6 & 486 & 57, 58 \\
R2 & 41--69, 76--81, 116 & 16 & 1766 & 57, 59, 60, 61 \\
R3 & 94--95, 98--112 & 9 & 509 & 58, 59 \\
R4 & 17, 21--32, 113--115 & 7 & 305 & 54, 55, 62, 63, 64 \\
R5 & 15, 18--20, 33--44 & 6 & 486 & 55, 60, 62 \\
R6 & 70--75, 118 & 4 & 232 & 61, 63 \\
R7 & 8--10 & 2 & 28 & 56 \\
\bottomrule
\end{tabular}
\end{table}

\section{Paired Evaluation Protocol and Parameters}\label{app:scenario_experiments}
All six pilots follow the MACPO outer loop in Section~\ref{sec:experiments}: LLSO local search, gated or always-on negotiation on shared variables, stop at rank-synchronized evaluation cap. Both methods share matched random seeds, swarm size 100, and penalty weight $w=|f|/512$. Proposed uses the complete configuration with fail-safe period $K{=}2$ on IEEE cases.

Improvement metric:
\begin{equation}
\Delta=\frac{F^{\mathrm{MACPO}}-F^{\mathrm{Proposed}}}{\bigl|F^{\mathrm{MACPO}}\bigr|}\times 100\%,
\label{eq:app_delta}
\end{equation}
where $\Delta>0$ favors Proposed (lower terminal $F$; more negative is better in EV dispatch). IEEE cases additionally report median paired $\Delta$, Wilcoxon $p$, and success (Proposed $F\le 2\times$ paired MACPO).

\begin{algorithm}[H]
\caption{Paired application-case evaluation}
\label{alg:app_paired}
\begin{algorithmic}[1]
\For{run $k=1$ to $25$}
  \State draw paired seed $s$
  \State run always-on MACPO ($g_i^t\equiv 1$) with LLSO until evaluation cap
  \State run complete Proposed configuration with the same seed and cap
  \State record best $F$, $\bar p_{\mathrm{comm}}$, and outer iterations
\EndFor
\end{algorithmic}
\end{algorithm}

\begin{algorithm}[H]
\caption{Paired IEEE dispatch evaluation}
\label{alg:ieee_paired}
\begin{algorithmic}[1]
\For{run $k=1$ to $25$}
  \State draw paired seed $s$; set fail-safe $K{=}2$
  \State run always-on MACPO with $n_{\mathrm{reg}}$ MPI processes until evaluation cap
  \State run complete Proposed configuration with the same seed and cap
  \State record $F$, $\bar p_{\mathrm{comm}}$; compute per-run $\Delta_k$; flag success
\EndFor
\State aggregate median $\Delta$, mean $\pm$ std, Wilcoxon $p$
\end{algorithmic}
\end{algorithm}

\begin{table}[H]
\centering
\caption{Application-case settings (25 paired runs each). Cap: maximum function evaluations per run ($d_{\max}$: maximum local dimension across agents).}
\label{tab:application_setup}
\begin{tabular}{@{}lcccc@{}}
\toprule
\textbf{Scenario} & \textbf{Agents} & \textbf{$D$} & \textbf{Gen./dim.} & \textbf{Eval.\ cap} \\
\midrule
Resource & 10 & 19 & 3.0 & $\approx 6.0\times 10^4$ \\
EV dispatch & 10 & 19 & 3.0 & $\approx 6.0\times 10^4$ \\
MAED-13 & 1 & 13 & 4.0 & $3.25\times 10^5$ \\
IEEE-30 & 4 & 11 & 2.0 & $\approx 2.2\times 10^4$ \\
IEEE-57 & 4 & 11 & 2.0 & $\approx 2.6\times 10^4$ \\
IEEE-118 & 8 & 65 & 2.5 & $\approx 1.7\times 10^5$ \\
\bottomrule
\end{tabular}
\end{table}

\section{Run-Level Statistics and Result Analysis}\label{app:applications}
Table~\ref{tab:application_run_stats} gives run-level $\Delta$ quartiles for all six pilots. On resource scheduling, Proposed improves mean terminal $F$ by 11.0\% while cutting communication by 88.4\%; median $\Delta{=}+9.5\%$ with one MACPO outlier widening variance. EV dispatch improves mean net benefit by 10.3\% (median $+0.6\%$). MAED-13 changes mean cost by only $-0.05\%$ with 88.7\% communication drop; all runs have $|\Delta|<1.5\%$. On IEEE~30/57, median paired $\Delta$ is near zero with 41--42\% communication reduction; IEEE~118 trades $-0.6\%$ median cost ($p{=}0.01$) for 60\% fewer negotiations. IEEE~30 success is 23/25 due to two LLSO stagnation outliers, not under-triggering.

\begin{table}[H]
\centering
\caption{Run-level paired $\Delta$ statistics ($\Delta=(\mathrm{MACPO}-\mathrm{Proposed})/|\mathrm{MACPO}|\times100\%$; positive favors Proposed).}
\label{tab:application_run_stats}
\begin{tabular}{@{}lrrr@{}}
\toprule
Scenario & Median $\Delta$ (\%) & Q1--Q3 (\%) & Min--Max (\%) \\
\midrule
\multicolumn{4}{@{}l@{}}{\textit{MACPO-style simulators} (RL comm.\ 11.3--11.6\%)} \\
MAED-13 & 0.0 & [0.0, 0.0] & [$-$1.5, +0.4] \\
Resource scheduling & +9.5 & [$-$0.1, +19.5] & [$-$25.8, +42.9] \\
EV dispatch & +0.6 & [$-$1.2, +3.4] & [$-$4.2, +88.7] \\
\midrule
\multicolumn{4}{@{}l@{}}{\textit{IEEE regional dispatch} (RL comm.\ 40.0--58.8\%)} \\
IEEE-30 & +0.5 & [$-$0.3, +1.2] & [$-$2.1, +4.8] \\
IEEE-57 & $-$0.0 & [$-$0.4, +0.3] & [$-$1.8, +1.5] \\
IEEE-118 & $-$0.6 & [$-$0.9, $-$0.3] & [$-$1.4, $-$0.1] \\
\bottomrule
\end{tabular}
\end{table}
